%% Copyright 2007, 2008, 2009 Elsevier Ltd
%%
%% This file is part of the 'Elsarticle Bundle'.
%% ---------------------------------------------
%%
%% It may be distributed under the conditions of the LaTeX Project Public
%% License, either version 1.2 of this license or (at your option) any
%% later version.  The latest version of this license is in
%%    http://www.latex-project.org/lppl.txt
%% and version 1.2 or later is part of all distributions of LaTeX
%% version 1999/12/01 or later.
%%
%% The list of all files belonging to the 'Elsarticle Bundle' is
%% given in the file `manifest.txt'.
%%

%% Template article for Elsevier's document class `elsarticle'
%% with numbered style bibliographic references
%% SP 2008/03/01
%%
%%
%%
%% $Id: elsarticle-template-num.tex 4 2009-10-24 08:22:58Z rishi $
%%
%%
%% \documentclass[preprint,5p,onecolumn,authoryear]{elsarticle}
\documentclass[authoryear,preprint,review,10pt]{elsarticle}
%% Use the option review to obtain double line spacing
%% \documentclass[preprint,review,12pt]{elsarticle}

%% Use the options 1p,twocolumn; 3p; 3p,twocolumn; 5p; or 5p,twocolumn
%% for a journal layout:
%% \documentclass[final,1p,times]{elsarticle}
%% \documentclass[final,1p,times,twocolumn]{elsarticle}
%% \documentclass[final,3p,times]{elsarticle}

%% \documentclass[authoryear,final,3p,times,twocolumn]{elsarticle}

%% \documentclass[final,5p,times]{elsarticle}
%% \documentclass[final,5p,times,twocolumn]{elsarticle}

%% if you use PostScript figures in your article
%% use the graphics package for simple commands
%% \usepackage{graphics}
%% or use the graphicx package for more complicated commands
%% \usepackage{graphicx}
%% or use the epsfig package if you prefer to use the old commands
%% \usepackage{epsfig}

%% The amssymb package provides various useful mathematical symbols
\usepackage{amssymb}
\usepackage{float}
\usepackage{eurosym} % para el euro
%% The amsthm package provides extended theorem environments
%% \usepackage{amsthm}
\usepackage[hyphens]{url}
\usepackage{hyperref}
\hypersetup{colorlinks,
%citecolor=blue
}
\usepackage{rotating}
\usepackage{ulem} % Paquete para tachar texto
\usepackage{geometry}
\usepackage{adjustbox}
\geometry{margin=1in}
\usepackage{longtable}
\usepackage{framed} % Framing content
\usepackage{multicol} % Multiple columns environment
\usepackage{nomencl} % Nomenclature package
\usepackage{mathtools}
\usepackage{cuted}
\makenomenclature
\setlength{\nomitemsep}{-\parskip} % Baseline skip between items
\DeclareMathOperator*{\argmax}{argmax}
\usepackage{array,multirow,booktabs}
\usepackage{makecell}
\renewcommand{\cellalign}{bl}
\usepackage{siunitx}
\usepackage{threeparttable}

\usepackage{breqn}
\usepackage{amsmath}
%\usepackage{showframe}
\usepackage{threeparttable}
%\usepackage{setspace}
\DeclarePairedDelimiter\ceil{\lceil}{\rceil}
\DeclarePairedDelimiter\floor{\lfloor}{\rfloor}    
    
\usepackage{natbib}
\usepackage{soul,color}
\soulregister\cite7
\soulregister\ref7
\soulregister\pageref7
\usepackage{textcomp}
\usepackage{gensymb}
\usepackage{caption}
\usepackage{boldline} 
\captionsetup[figure]{name={\textbf{Fig.}}, labelfont=bf, justification=centering}
\captionsetup[table]{justification=raggedright,  labelsep=newline, singlelinecheck=false, labelfont=bf}
\usepackage{subcaption}
\usepackage[ruled,vlined]{algorithm2e}
\usepackage{lscape} 
\usepackage{longtable} 

\renewcommand*\nompreamble{\begin{multicols}{2}}
\renewcommand*\nompostamble{\end{multicols}}
    
%% The lineno packages adds line numbers. Start line numbering with
%% \begin{linenumbers}, end it with \end{linenumbers}. Or switch it on
%% for the whole article with \linenumbers after \end{frontmatter}.
%% \usepackage{lineno}

%% natbib.sty is loaded by default. However, natbib options can be
%% provided with \biboptions{...} command. Following options are
%% valid:

%%   round  -  round parentheses are used (default)
%%   square -  square brackets are used   [option]
%%   curly  -  curly braces are used      {option}
%%   angle  -  angle brackets are used    <option>
%%   semicolon  -  multiple citations separated by semi-colon
%%   colon  - same as semicolon, an earlier confusion
%%   comma  -  separated by comma
%%   numbers-  selects numerical citations
%%   super  -  numerical citations as superscripts
%%   sort   -  sorts multiple citations according to order in ref. list
%%   sort&compress   -  like sort, but also compresses numerical citations
%%   compress - compresses without sorting
%%
%% \biboptions{comma,round}

% \biboptions{}
\usepackage{ulem}

\journal{Utilities Policy}

\begin{document}

\begin{frontmatter}

\title{Evaluating Residential Cooking Electrification Policies under Tariff Uncertainty: A Case Study of Bogotá, Colombia}

\author[label1]{Julio C. Cruz\corref{cor1}}
\cortext[cor1]{Corresponding author}
\address[label1]{Department of Industrial Engineering, University of Los Andes,  Cra 1 Este No 19A-40, Bogot\'a, Colombia}
\ead{jc.cruz412@uniandes.edu.co}

\author[label1]{Sergio Cabrales}
\ead{s-cabral@uniandes.edu.co}

\begin{abstract}
When households switch from natural gas to induction cooking, they shift energy use between two regulated services with different tariffs, subsidies, and customer contributions. This study examines how those rules shape household costs and the net fiscal impact of cooking electrification in Bogotá, Colombia, under tariff uncertainty. We reconstruct regulated gas and electricity tariffs, simulate their drivers using Monte Carlo methods, and account for household demand responses to prices. We then estimate ten-year total ownership costs, changes in subsidies and contributions, greenhouse gas emissions, and marginal abatement costs. Under the modeled conditions, induction lowers expected ownership costs across all six tariff strata, with ten-year savings ranging from USD~253 per household in Stratum~1 to USD~2,532 in Stratum~6. The modeled switch yields expected aggregate net fiscal savings of USD~538 million, mainly because avoided gas subsidies outweigh the loss of contribution revenues. Greenhouse gas emissions decline across strata, and the estimated intervals for marginal abatement costs remain below zero. Scenarios with 50\% and 100\% imported gas increase the estimated household and fiscal savings. The results highlight two priorities for utility policy: assessing upfront appliance and installation costs when designing targeted support, and tracking changes in both subsidies and contributions as cooking demand shifts from gas to electricity.
\end{abstract}

%% keywords here, in the form: keyword \sep keyword

%% MSC codes here, in the form: \MSC code \sep code
%% or \MSC[2008] code \sep code (2000 is the default)
\begin{keyword}
Cooking electrification \sep Tariff regulation \sep Cross-subsidies \sep Energy affordability \sep Tariff uncertainty
\end{keyword}

\end{frontmatter}

%%
%% Start line numbering here if you want
%%
% \linenumbers

%% main text

%%%%%%%%%%%%%%%%%%%%%%%%%%%%%%%%%%%%%%%%%%%%%%%%%%%%%%%
\section{Introduction}
\label{introduction}

Replacing a gas stove with an induction cooktop shifts household energy use from the gas network to the electricity network. In regulated markets, the two services can differ in how costs are passed through to customers and how subsidies and contributions are assigned. These differences affect both the savings households can expect and the net cost of supporting residential energy consumption. Current gas and electricity prices provide a snapshot of that trade-off, while the decision to replace a stove depends on an upfront investment and energy bills over the appliance's lifetime.

Cooking electrification can also reduce greenhouse gas (GHG) emissions, depending on the electricity mix, appliance efficiency, and household use patterns \citep{yangka2016,im2020,lee2021,bandekar2022}. Total cost of ownership (TCO) and marginal abatement cost curves (MACCs) help assess whether these emissions reductions come with additional costs or savings for households \citep{cabrales2023costos,cimini2022,floess2023}. An assessment based on fixed tariffs, however, gives limited insight into how future tariff changes affect the size and distribution of household savings. It also leaves open a separate question for policymakers: how the shift in demand changes the subsidies paid and the contributions collected across the two services.

Bogotá provides a case in which these effects can be examined across six tariff strata. In Colombia, regulated gas and electricity tariffs reflect sector-specific cost formulas, macroeconomic indexation, and cross-subsidy rules that differentiate what households pay by socioeconomic stratum \citep{CREG119_2007,CREG137_2013,CREG183_2013,CREG184_2014,cabrales2023costos}. Concerns about domestic gas availability and greater reliance on imports add uncertainty to future gas prices \citep{ramirez2025macro_gas,cabrales2025politica_gas}. The study therefore also considers two gas-import scenarios, with 50\% and 100\% imported gas, using tariff increases from \citet{cabrales2025fedesarrollo_gas_imports}.

This study asks how regulated gas and electricity tariffs and their cross-subsidy rules shape household savings across tariff strata and the net fiscal effect of replacing gas stoves with induction cooktops under tariff uncertainty. To answer this question, we reconstruct regulated tariffs, simulate their underlying drivers, and account for price-responsive household demand. We compare ten-year ownership costs by stratum, estimate changes in subsidies and contributions, and quantify GHG emissions and marginal abatement costs. Bringing these outcomes together shows how the same switch can affect household affordability and the financing of cross-subsidies. The analysis informs the design of targeted support for appliance acquisition and household adaptations, and identifies which subsidy and contribution flows policymakers should monitor as cooking demand shifts between services.

%%%%%%%%%%%%%%%%%%%%%%%%%%%%%%%%%%%%%%%%%%%%%%%%%%%%%%%
\section{Literature review}
\label{literature}

In regulated energy markets, tariff design shapes both the cost of cooking and the distribution of subsidies and contributions. Switching from gas to electricity changes the tariff a household faces and can alter its subsidy or contribution payments across the two services. In Colombia, \citet{cabrales2023costos} estimate household savings from cooking electrification across cities and socioeconomic strata while also examining national investment requirements and fiscal effects. Under some of their national scenarios, lower gas subsidies are offset by higher electricity-related fiscal costs. Complementary work examines the distributional and emissions implications of gas subsidies \citep{cabrales2025subsidios_gas}. These studies show why household savings and the net fiscal impact need to be assessed separately.

Household savings depend on the purchase price of the appliance and on energy expenditures over its lifetime \citep{dagnachew2020scenario}. Total cost of ownership (TCO) combines acquisition and operating costs, allowing a comparison between appliances with different upfront costs and energy efficiencies \citep{ellram1994taxonomy}. Marginal abatement cost curves (MACCs) relate the difference in costs to the associated emissions reductions. Evidence from Bhutan and Korea, together with comparative and life-cycle studies, shows that these reductions depend on electricity carbon intensity, appliance efficiency, and use patterns \citep{yangka2016,im2020,lee2021,bandekar2022,cimini2022}. Evidence from developing regions also highlights the role of infrastructure and household conditions in shaping the transition \citep{dagnachew2020scenario,floess2023,scott2024energy}.

A comparison at fixed tariffs provides a benchmark for household savings. To assess how those savings vary under uncertainty, changes in energy costs and macroeconomic indices must be translated into the tariffs households face. Research on energy prices and macroeconomic variables provides methods for modeling these drivers \citep{su2019data,tan2010day,khamis2020,rosales2017,babu2015}. Tariff changes can also affect consumption. The meta-analysis by \citet{espey2004turning} finds that residential electricity demand falls as prices rise, with stronger responses over longer horizons. Applying residential demand elasticities to cooking requires an assumption about how consumption for this specific end use responds to prices.

The question addressed here is how uncertainty in regulated gas and electricity tariffs affects household savings and the balance of subsidies and contributions when cooking shifts between the two services. This study links simulated tariff components to the rates paid by each socioeconomic stratum in Bogotá and carries that uncertainty through to household costs, fiscal effects, emissions, and marginal abatement costs. Evaluating these outcomes together shows how a change that reduces household costs can also alter the financing of residential energy subsidies.


%%%%%%%%%%%%%%%%%%%%%%%%%%%%%%%%%%%%%%%%%%%%%%%%%%%%%%%%%%
\section{Methodology}
\label{methodology}

The framework compares continued natural gas cooking under business as usual (BAU) with replacement by induction cooktops. It estimates household total cost of ownership (TCO), the net fiscal effect of changes in subsidies and contributions, greenhouse gas (GHG) emissions, and marginal abatement costs (MAC) \citep{ellram1994taxonomy,clerc2013desarrollo}. The analysis combines a deterministic benchmark with Monte Carlo simulations to assess how uncertainty in regulated tariffs affects these outcomes across socioeconomic strata.

Figure~\ref{fig:methodology_diagram} summarizes the six stages of the analysis: baseline construction, tariff and demand simulation, parameter selection, ownership-cost and fiscal analysis, emissions estimation, and construction of marginal abatement cost curves (MACCs).


\begin{figure}[H]
    \centering
    \includegraphics[width=\textwidth]{Methodology_Electrification_Flowchart2.jpg}
    \caption{Methodological framework for evaluating residential cooking electrification under regulated tariffs and price-responsive household demand.}
    \label{fig:methodology_diagram}
\end{figure}

%------------------------------------------------------------------------
\subsection{Regulated tariff structure}
\label{sec:regulated_tariffs}
%------------------------------------------------------------------------

Colombia's Comisión de Regulación de Energía y Gas (CREG) defines the components used to calculate residential gas and electricity unit costs. These costs provide a reference for determining the tariffs households face after applying the subsidy and contribution rules for each socioeconomic stratum.

The reference expression for the volumetric unit cost of residential natural gas is

\begin{equation}
Cuv_t = \frac{G_t + T_t}{1-\rho_t} + \left(D_t + fpc_t\right),
\label{eq:gas_tariff_method}
\end{equation}
 \citep{CREG137_2013,CREG183_2013,CREG184_2014}

where $G_t$ is the gas supply component, $T_t$ is the transportation component, $D_t$ is the distribution component, $fpc_t$ is the term associated with the calorific adjustment, and $\rho_t$ represents technical losses.

The reference expression for the residential electricity unit cost is


\begin{equation}
CU_t = GM_t + T_t + D_{n,t} + C_{v,t} + PR_{n,t} + R_t,
\label{eq:elec_tariff_method}
\end{equation}
\citep{CREG119_2007}

where $GM_t$ is the generation component, $T_t$ is the transmission charge, $D_{n,t}$ is the distribution component, $C_{v,t}$ is the commercialization margin, $PR_{n,t}$ represents recognized losses, and $R_t$ captures system restriction costs. Stage~2 describes how the simulation uses these cost components to construct monthly tariff paths.

%------------------------------------------------------------------------
\subsection{Stage 1. Data collection and baseline construction}
\label{sec:data_collection}
%------------------------------------------------------------------------

The baseline combines residential natural gas consumption, gas and electricity tariffs and their components, subscriber counts by socioeconomic stratum, and macroeconomic series used for tariff indexation. The data are obtained from official regulatory and statistical sources. Gas consumption is expressed in cubic metres per month and electricity consumption in kilowatt-hours per month. Prices and quantities are expressed in compatible units for calculating household expenditure. Section~\ref{case study} describes the data sources and baseline for Bogotá.

%------------------------------------------------------------------------
\subsection{Stage 2. Stochastic simulation of tariff drivers and price-responsive demand}
\label{sec:variable_simulation}
%------------------------------------------------------------------------

The simulation focuses on three drivers of residential tariffs: the natural gas supply component ($G_t$), the electricity generation component ($GM_t$), and inflation ($\pi_t$). The first two enter the respective unit costs directly, while inflation affects the components subject to indexation. Model specification, selection, and validation draw on established approaches to energy-price modeling \citep{lucia2002electricity,tan2010day,ibrahim2021geometric}. Appendix~\ref{appendix:stochastic_simulation} provides the technical details.

For each Monte Carlo path $i$, the simulation combines the monthly gas supply component with a block of other gas tariff components updated using simulated inflation. The electricity calculation combines the generation component with the remaining electricity tariff components. Subsidies and contributions are then applied to obtain the prices used to calculate household expenditure for each stratum.

The demand module uses a price-elasticity adjustment informed by evidence on residential energy demand \citep{espey2004turning}. The adjustment produces monthly consumption paths that feed into the calculations of household expenditure, fiscal effects, and emissions. Appendix~\ref{appendix:demand_adjustment} describes how the adjustment is applied.

The base case assumes no imported gas. The sensitivity analysis considers gas import shares of 50\% and 100\%, represented by increases of 44.6\% and 89.3\%, respectively, in the corresponding base-case gas tariff in each simulated month \citep{cabrales2025fedesarrollo_gas_imports}. Electricity tariffs and consumption paths are held fixed across the three import scenarios. This comparison isolates the effect of the gas tariff shock under the modeled conditions.

%------------------------------------------------------------------------
\subsection{Stage 3. Technical and financial parameters}
\label{sec:parameter_selection}
%------------------------------------------------------------------------

The analysis specifies the evaluation horizon, monthly discount rate, appliance capital cost, appliance efficiencies, emission factors, and price elasticity. The two technologies are compared using a common useful cooking-energy requirement. Their efficiencies determine the final energy needed to deliver that service. Consumption is fixed by stratum in the deterministic evaluation and varies across simulation paths in the stochastic evaluation. Section~\ref{case study} reports the parameter values used for Bogotá.

%------------------------------------------------------------------------
\subsection{Stage 4. Total cost of ownership and fiscal effects}
\label{sec:cost_calculation}
%------------------------------------------------------------------------

For each socioeconomic stratum and simulation path $i$, the TCO of cooking option $x \in \{\text{gas},\text{electricity}\}$ is
\begin{equation}
TCO_x^{(i)} =
CAPEX_x + \sum_{t=1}^{T}
\frac{OPEX_{x,t}^{(i)}}{(1+r)^t},
\end{equation}
where $CAPEX_x$ is the upfront appliance cost, $r$ is the monthly discount rate, and $T$ is the number of months in the evaluation horizon. This formulation combines the initial investment with the present value of operating expenditure \citep{ellram1994taxonomy,bauer2025total,ElectricVehicleTCO}.

Monthly operating expenditure is
\begin{equation}
OPEX_{x,t}^{(i)} =
Price_{x,t}^{(i)} \times Consumption_{x,t}^{(i)},
\end{equation}
where $Price_{x,t}^{(i)}$ is the tariff paid by the household after applying the relevant subsidy or contribution, and $Consumption_{x,t}^{(i)}$ is the corresponding final energy used for cooking. The change in ownership costs is
\begin{equation}
\Delta TCO^{(i)} =
TCO_{\text{electricity}}^{(i)} - TCO_{\text{gas}}^{(i)}.
\end{equation}
A negative value indicates household savings relative to continued gas cooking. The Monte Carlo paths yield a distribution of this change for each stratum.

For the fiscal analysis, we calculate the monthly change in subsidy expenditure net of contribution revenue when cooking shifts from gas to induction. For each stratum and simulation path, we discount 120 monthly flows at a rate of 0.2919\% per month. We multiply the resulting household-level present value by Bogotá's December 2022 residential gas subscriber count for the corresponding stratum. The citywide fiscal impact is obtained by summing the stratum-level effects within each path. A negative net fiscal change represents a saving under the modeled replacement scenario.

%------------------------------------------------------------------------
\subsection{Stage 5. GHG emissions}
\label{sec:ghg_calculation}
%------------------------------------------------------------------------

Monthly GHG emissions are calculated by multiplying final energy use for cooking by the corresponding emission factor. For each cooking option $x$ and simulation path $i$, monthly emissions and avoided emissions are
\begin{align}
GHG_{x,t}^{(i)}
    &= EF_x \times Consumption_{x,t}^{(i)}, \\
\Delta GHG_t^{(i)}
    &= GHG_{\text{gas},t}^{(i)}
       - GHG_{\text{electricity},t}^{(i)},
\end{align}
where $EF_x$ is the emission factor for energy carrier $x$, expressed in units compatible with final energy consumption. A positive $\Delta GHG_t^{(i)}$ indicates that induction reduces emissions relative to BAU. Emission factors are held fixed, while consumption varies across simulation paths.

For the abatement cost calculation, monthly emissions reductions are discounted over the evaluation horizon using the same monthly rate as the ownership-cost calculation. With emissions expressed in tonnes of CO$_2$e, discounted avoided emissions are

\begin{equation}
\Delta GHG_{\mathrm{disc}}^{(i)}
    = \sum_{t=1}^{T}
      \frac{\Delta GHG_t^{(i)}}{(1+r)^t}.
\end{equation}

%------------------------------------------------------------------------
\subsection{Stage 6. Marginal abatement cost and MACC}
\label{sec:mac}
%------------------------------------------------------------------------

The MAC relates the change in household ownership costs to discounted avoided emissions. Following the approach to abatement cost analysis under uncertainty discussed by \citet{clerc2013desarrollo}, we calculate

\begin{equation}
MAC^{(i)} =
\frac{\Delta TCO^{(i)}}{\Delta GHG_{\mathrm{disc}}^{(i)}}.
\end{equation}

When avoided emissions are positive, a negative MAC indicates that electrification reduces emissions and saves households money. A positive MAC indicates that the emissions reduction comes at an additional household cost. Calculating the ratio across simulation paths yields a distribution of MAC values for each stratum.

The MACC orders the strata from the lowest to the highest cost per tonne of CO$_2$e avoided \citep{clerc2013desarrollo}. The height of each bar represents the stratum's marginal abatement cost, and its width represents the associated emissions reduction on the abatement basis used for the curve. The horizontal axis accumulates those reductions across strata.

%%%%%%%%%%%%%%%%%%%%%%%%%%%%%%%%%%%%%%%%%%%%%%%%%%%%%%%
\section{Case Study}
\label{case study}

Bogotá has a large residential natural gas customer base and a tariff structure that differentiates prices across six socioeconomic strata. It therefore provides a case for examining how a shift from gas to induction cooking affects household costs, emissions, and the balance of subsidies and contributions across the two regulated services.

Stratum~4 serves as the reference because its residential tariffs include neither a subsidy nor a contribution. The analysis uses this benchmark to characterize the underlying regulated costs and then evaluates the switch for all six strata, accounting for differences in consumption and the tariff adjustments applicable to each service.

%------------------------------------------------------------------------
\subsection{Context and data sources}
\label{sec:context_data_sources}
%------------------------------------------------------------------------

The analysis draws primarily on administrative records from Colombia's Sistema Único de Información de Servicios Públicos Domiciliarios (SUI). These records provide residential tariffs and tariff components, gas consumption, and subscriber counts by stratum. The December~2022 snapshot provides the subscriber counts used for citywide aggregation, while the monthly series provide the consumption and tariff histories used in the analysis. Appendix~\ref{appendix:case_context} describes the national gas-market context and the composition of Bogotá's residential gas customer base.

%------------------------------------------------------------------------
\subsection{Data description}
%------------------------------------------------------------------------
\subsubsection{Macroeconomic variables}
%------------------------------------------------------------------------

The macroeconomic series include the Consumer Price Index (IPC), the Producer Price Index (IPP), and the representative market exchange rate (TRM). The IPC and IPP provide context for the indexation of regulated cost components, while the TRM is used to express monetary series in U.S.\ dollars. The stochastic model uses inflation ($\pi_t$) alongside the natural gas supply component ($G_t$) and the electricity generation component ($GM_t$), as described in Section~\ref{sec:variable_simulation}.

%------------------------------------------------------------------------
\subsubsection{Residential gas subscribers and consumption}
%------------------------------------------------------------------------

SUI records identify the number of residential gas subscribers in each stratum. These counts are held fixed in the citywide aggregation. Strata~1--3 are eligible for subsidies under the applicable rules for each service, Strata~5 and~6 pay contributions, and Stratum~4 serves as the reference \citep{CREG137_2013,CREG184_2014,cabrales2023costos}. The resulting differences in end-user tariffs affect household expenditures and determine how a shift in consumption between gas and electricity changes subsidy expenditure and contribution revenue.

Monthly gas consumption per subscriber provides the starting point for the household analysis. Figure~\ref{fig:cons-by-stratum} compares Bogotá with the other cities among the ten largest residential gas markets and with the subscriber-weighted average for those ten cities. Differences in consumption across strata affect the magnitude of the costs and emissions associated with the modeled switch.

\begin{figure}[H]
    \centering
    \includegraphics[width=0.98\linewidth]{NatGasConsbyStratum.png}
    \caption{Monthly residential gas consumption per subscriber by city and socioeconomic stratum (m$^3$/month), with the subscriber-weighted average for the ten Colombian cities with the largest residential gas customer bases in December~2022.}
    \label{fig:cons-by-stratum}
\end{figure}

%------------------------------------------------------------------------
\subsubsection{Regulated gas and electricity tariffs}
%------------------------------------------------------------------------

The gas tariff analysis uses the volumetric unit cost ($Cuv_t$) described in Eq.~\eqref{eq:gas_tariff_method}. The gas supply component reflects procurement conditions, while transportation and distribution charges follow regulated cost and adjustment rules \citep{CREG137_2013,CREG183_2013,CREG184_2014,CREG008_2014}. Figure~\ref{fig:price-gas-bog-e4} shows the evolution of the gas unit cost for Stratum~4.

For electricity, the analysis uses the regulated unit cost ($CU_t$) described in Eq.~\eqref{eq:elec_tariff_method}. The generation component reflects the cost of electricity purchases through contracts and the wholesale market, while the remaining components account for network charges, commercialization, recognized losses, and system restrictions \citep{CREG119_2007}. Figure~\ref{fig:price-elec-bog-e4} shows the corresponding series for Stratum~4. Both histories exhibit changes over time and provide the tariff context for the stochastic analysis.

%------------------------------------------------------------------------
\subsubsection{Relative cost of useful cooking energy}
%------------------------------------------------------------------------

The gas and electricity series are converted to a common energy basis for comparison. Figure~\ref{fig:energy-cost-comparison-corr} reports the cost per unit of \emph{useful} cooking energy by dividing each tariff, expressed in USD/MMBtu of final energy, by the corresponding stove efficiency. This comparison shows how relative tariffs and appliance efficiencies jointly determine cooking energy costs. The total cost of ownership calculation also includes the upfront investment in induction. Appendix~\ref{appendix:case_context} reports the comparison on a final-energy basis.

\begin{figure}[H]
    \centering
    \includegraphics[width=0.95\linewidth]{ComparaGasElec.png}
    \caption{Cost of residential cooking energy in Bogotá (USD/MMBtu of useful energy). The natural gas and electricity tariffs are adjusted using the assumed efficiencies of gas stoves and induction cooktops.}
    \label{fig:energy-cost-comparison-corr}
\end{figure}

%------------------------------------------------------------------------
\subsubsection{Technical and financial parameters}
%------------------------------------------------------------------------

Table~\ref{tab:parameters} summarizes the assumptions used to compare continued gas cooking with induction. The analysis uses a ten-year horizon.

\begin{table}[H]
    \centering
    \caption{Technical and financial parameters for the Bogotá case study}
    \label{tab:parameters}
    \begin{tabular}{p{0.22\linewidth} p{0.34\linewidth} p{0.34\linewidth}}
        \hline
        \textbf{Category} & \textbf{Parameter} & \textbf{Value / Source} \\
        \hline
        Investment
        & Induction cooking system
        & USD~316.9 \citep{upme2018lena} \\
        \hline
        Evaluation
        & Analysis horizon ($T$)
        & 120 months (10 years) \\
        & Monthly discount rate ($r$)
        & 0.2919\% (equivalent to 3.56\% annually, based on TNX) \\
        \hline
        Appliance efficiency
        & Induction cooktop
        & 90\% \citep{sadhu2010induction} \\
        & Gas stove
        & 40\% \citep{doe2008cooking} \\
        \hline
        Demand response
        & Assumed short-run price elasticity ($\varepsilon_x$)
        & $-0.35$, informed by \citet{espey2004turning} \\
        \hline
    \end{tabular}
\end{table}


%%%%%%%%%%%%%%%%%%%%%%%%%%%%%%%%%%%%%%%%%%%%%%%%%%%%%%%
\section{Results}
\label{results}

The Bogotá case study compares continued cooking with natural gas under business as usual (BAU) with replacement by induction cooktops. The deterministic evaluation uses historical tariffs, while the stochastic evaluation propagates simulated tariff and demand paths through the household cost, fiscal, and emissions calculations described in Section~\ref{methodology}. The results first cover the base case with no imported gas, followed by the 50\% and 100\% gas import scenarios.

%------------------------------------------------------------------------
\subsection{Stochastic simulation of tariff drivers}
\label{results:simulation}
%------------------------------------------------------------------------

The selected specification models the natural gas supply component ($G_t$) with heavy-tailed innovations and allows the electricity generation component ($GM_t$) to depend on the gas supply component. This dependence matters because changes in gas supply costs can affect the tariffs faced under both cooking options. \ref{appendix:stochastic_diagnostics} reports the dependence diagnostics, model selection, backtesting, and simulated trajectories.

%------------------------------------------------------------------------
\subsection{Total cost of ownership}
\label{results:tco}
%------------------------------------------------------------------------

Table~\ref{tab:tco_ci} reports the ten-year ownership cost of continued gas cooking in the deterministic benchmark and the change in cost when households switch to induction. The change, $\Delta TCO$, includes the upfront cost of the induction cooktop. Negative values indicate savings relative to BAU.

\begin{table}[htbp]
    \centering
    \caption{Ten-year household ownership costs and changes under electrification by socioeconomic stratum in Bogotá (USD per household, present value).}
    \label{tab:tco_ci}
    \begin{tabular}{lrrc}
        \toprule
        & \textbf{BAU gas TCO} & \multicolumn{2}{c}{\textbf{Change under electrification ($\Delta TCO$)}} \\
        \cmidrule(lr){3-4}
        \textbf{Stratum} & \textbf{Deterministic} & \textbf{Deterministic}
        & \textbf{Stochastic mean [95\% interval]} \\
        \midrule
        E1 & 680.11  & -287.47  & $-252.57\;[-357.68,-149.78]$ \\
        E2 & 710.50  & -310.94  & $-304.51\;[-419.20,-192.34]$ \\
        E3 & 955.86  & -531.83  & $-709.01\;[-892.41,-530.30]$ \\
        E4 & 1177.21 & -705.08  & $-918.61\;[-1146.65,-695.59]$ \\
        E5 & 1777.17 & -1225.92 & $-1548.27\;[-1892.53,-1211.59]$ \\
        E6 & 3212.43 & -2537.55 & $-2532.08\;[-3057.92,-2017.81]$ \\
        \bottomrule
    \end{tabular}
\end{table}

Switching to induction lowers estimated ownership costs in all six strata. Mean savings under the stochastic evaluation range from USD~252.57 per household in E1 to USD~2,532.08 in E6. The intervals widen in absolute dollar terms at higher strata. Even at the upper bound of each reported interval, the present value of lower operating expenditures exceeds the upfront appliance cost.

%------------------------------------------------------------------------
\subsection{Fiscal impact of electrification}
\label{results:fiscal}
%------------------------------------------------------------------------

The fiscal analysis measures the change in subsidy expenditure net of contribution revenue as cooking demand shifts from gas to electricity. Table~\ref{tab:fiscal_impact_aggregate} reports the ten-year present values for Bogotá, aggregated using the December~2022 residential gas subscriber counts by stratum. The BAU column reports the deterministic level of net expenditure under continued gas cooking. The remaining columns report the change relative to the corresponding BAU case. \ref{appendix:fiscalMACCandimportedgas} provides the per-household results.

\begin{table}[htbp]
    \centering
    \caption{Net subsidy expenditure under BAU and its change under electrification in Bogotá (million USD, ten-year present value). Negative changes indicate net fiscal savings.}
    \label{tab:fiscal_impact_aggregate}
    \begin{tabular}{lrrrc}
        \toprule
        & \textbf{BAU level} & \multicolumn{3}{c}{\textbf{Change under electrification}} \\
        \cmidrule(lr){3-5}
        \textbf{Stratum} & \textbf{Deterministic} & \textbf{Deterministic}
        & \textbf{Stochastic mean} & \textbf{95\% interval} \\
        \midrule
        E1 & 117.71 & -101.98 & -151.86 & $[-171.63,-132.46]$ \\
        E2 & 348.69 & -299.14 & -464.18 & $[-524.62,-404.88]$ \\
        E3 & 0.00   & 14.37   & 17.27   & $[13.54,22.33]$ \\
        E4 & 0.00   & 0.00    & 0.00    & $[0.00,0.00]$ \\
        E5 & -26.99 & 23.77   & 28.72   & $[25.05,32.45]$ \\
        E6 & -62.92 & 59.68   & 31.80   & $[26.91,37.87]$ \\
        \midrule
        \textbf{Total} & \textbf{376.49} & \textbf{-303.31}
        & \textbf{-538.26} & $\boldsymbol{[-630.76,-444.69]}$ \\
        \bottomrule
    \end{tabular}
\end{table}

Net fiscal savings are USD~303.31 million in the deterministic evaluation and average USD~538.26 million in the stochastic evaluation, with a reported 95\% interval of USD~444.69--630.76 million in savings. Lower net subsidy expenditure in E1 and E2 outweighs the increase in E3 and the reduction in net contribution revenue from E5 and E6. E4 has no fiscal change because it receives neither a subsidy nor a contribution charge in the modeled tariff structure.

%------------------------------------------------------------------------
\subsection{GHG emissions and abatement potential}
\label{results:ghg}
%------------------------------------------------------------------------

Table~\ref{tab:ghg_results} compares annual household emissions under continued gas cooking and electrification. Annual emissions decline in all six strata. Estimated reductions range from 0.064 to 0.151~tCO$_2$e per household per year in the deterministic evaluation and from 0.041 to 0.109~tCO$_2$e in the stochastic means. E6 has the largest absolute reduction in both evaluations, consistent with its higher baseline consumption.

\begin{table}[htbp]
    \centering
    \caption{Annual household GHG emissions from cooking in Bogotá (tCO$_2$e per household per year).}
    \label{tab:ghg_results}
    \begin{tabular}{lcccccc}
        \toprule
        \textbf{Scenario} & \textbf{E1} & \textbf{E2} & \textbf{E3}
        & \textbf{E4} & \textbf{E5} & \textbf{E6} \\
        \midrule
        BAU gas cooking                 & 0.450 & 0.393 & 0.317 & 0.391 & 0.492 & 0.751 \\
        Electrification (deterministic) & 0.360 & 0.314 & 0.253 & 0.312 & 0.393 & 0.600 \\
        Electrification (stochastic mean)
                                         & 0.385 & 0.352 & 0.271 & 0.334 & 0.420 & 0.642 \\
        \bottomrule
    \end{tabular}
\end{table}

%------------------------------------------------------------------------
\subsection{Marginal abatement cost curves}
\label{results:macc}
%------------------------------------------------------------------------

Figure~\ref{fig:macc-mean} presents the marginal abatement cost curve based on the mean stochastic MAC for each stratum. All six bars are below zero, indicating that induction reduces emissions at a net saving to households under the modeled conditions. The ordering is the same in the deterministic, mean, and median curves. \ref{appendix:fiscalMACCandimportedgas} reports the deterministic and median representations.

\begin{figure}[H]
    \centering
    \includegraphics[width=0.88\textwidth]{MACCBogMeanElec.png}
    \caption{Marginal abatement cost curve for residential cooking electrification in Bogotá, using the mean stochastic MAC for each socioeconomic stratum.}
    \label{fig:macc-mean}
\end{figure}

%------------------------------------------------------------------------
\subsection{Marginal abatement costs under tariff uncertainty}
\label{results:mac}
%------------------------------------------------------------------------

Table~\ref{tab:macc_ci} summarizes MAC values by stratum. The denominator is discounted avoided emissions over the evaluation horizon, as defined in Section~\ref{sec:mac}. Deterministic MAC values range from $-93.72$ to $-496.06$ USD per discounted tonne of CO$_2$e avoided, while stochastic means range from $-52.15$ to $-510.00$. The reported 95\% intervals are below zero in all strata.

\begin{table}[htbp]
    \centering
    \caption{Marginal abatement costs by socioeconomic stratum in Bogotá (USD per discounted tonne of CO$_2$e avoided, 200 simulations per stratum).}
    \label{tab:macc_ci}
    \begin{tabular}{lrrc}
        \toprule
        \textbf{Stratum} & \textbf{Deterministic} & \textbf{Stochastic mean}
        & \textbf{95\% interval} \\
        \midrule
        E1 & -93.72  & -52.15  & $[-81.72,-28.36]$ \\
        E2 & -116.12 & -100.02 & $[-141.66,-66.52]$ \\
        E3 & -246.05 & -289.57 & $[-368.78,-222.37]$ \\
        E4 & -264.86 & -307.00 & $[-392.41,-244.55]$ \\
        E5 & -366.06 & -439.00 & $[-537.57,-354.79]$ \\
        E6 & -496.06 & -510.00 & $[-619.67,-422.53]$ \\
        \bottomrule
    \end{tabular}
\end{table}

Figure~\ref{fig:boxplot-mac} shows the dispersion around these estimates. E5 and E6 combine the largest mean savings per tonne avoided with the widest reported intervals.

\begin{figure}[H]
    \centering
    \includegraphics[width=0.88\textwidth]{MACCBogEstocast.png}
    \caption{Distribution of simulated marginal abatement costs by socioeconomic stratum in Bogotá.}
    \label{fig:boxplot-mac}
\end{figure}

%------------------------------------------------------------------------
\subsection{Imported-gas sensitivity scenarios}
\label{results:import_scenarios}
%------------------------------------------------------------------------

The 50\% and 100\% gas import scenarios increase the residential gas tariff by 44.6\% and 89.3\%, respectively, relative to the case with no imported gas \citep{cabrales2025fedesarrollo_gas_imports}. Electricity tariffs and consumption paths are held fixed across scenarios. The resulting differences therefore reflect the effect of higher gas tariffs under the assumptions of the sensitivity analysis.

In the deterministic evaluation, the 50\% import scenario increases ten-year household savings to USD~590.80 in E1 and USD~3,970.30 in E6. At 100\% imports, the corresponding savings reach USD~894.81 and USD~5,406.25. The stochastic mean MAC remains negative in every stratum. It ranges from $-771.74$ to $-176.90$ USD per discounted tonne of CO$_2$e avoided at 50\% imports and from $-1{,}054.91$ to $-272.15$ at 100\%. \ref{appendix:fiscalMACCandimportedgas} reports the full stratum-level TCO and MAC results.

\begin{table}[htbp]
    \centering
    \caption{Deterministic net subsidy expenditure under BAU and its change under electrification in the imported-gas scenarios (million USD, ten-year present value). Negative changes indicate net fiscal savings.}
    \label{tab:import_scenarios_fiscal_aggregate}
    \begin{tabular}{lrrrr}
        \toprule
        & \multicolumn{2}{c}{\textbf{50\% imported gas}}
        & \multicolumn{2}{c}{\textbf{100\% imported gas}} \\
        \cmidrule(lr){2-3}\cmidrule(lr){4-5}
        \textbf{Stratum} & \textbf{BAU level} & \textbf{Change}
        & \textbf{BAU level} & \textbf{Change} \\
        \midrule
        E1 & 170.24 & -155.65 & 222.87  & -209.43 \\
        E2 & 504.22 & -457.06 & 660.09  & -615.33 \\
        E3 & 0.00   & 14.36   & 0.00    & 14.36 \\
        E4 & 0.00   & 0.00    & 0.00    & 0.00 \\
        E5 & -39.03 & 35.98   & -51.10  & 48.22 \\
        E6 & -90.99 & 88.08   & -119.12 & 116.54 \\
        \midrule
        \textbf{Total} & \textbf{544.43} & \textbf{-474.29}
        & \textbf{712.73} & \textbf{-645.64} \\
        \bottomrule
    \end{tabular}
\end{table}

Table~\ref{tab:import_scenarios_fiscal_aggregate} shows that deterministic net fiscal savings rise from USD~303.31 million in the base case to USD~474.29 million at 50\% imports and USD~645.64 million at 100\% imports. Higher gas tariffs increase both the subsidies avoided in E1 and E2 and the contribution revenue forgone in E5 and E6. In both scenarios, the increase in avoided subsidies is larger.

%%%%%%%%%%%%%%%%%%%%%%%%%%%%%%%%%%%%%%%%%%%%%%%%%%%%%%%
\section{Discussion}
\label{discussion}
%%%%%%%%%%%%%%%%%%%%%%%%%%%%%%%%%%%%%%%%%%%%%%%%%%%%%%%

The household cost results are consistent with the savings estimated for Colombia by \citet{cabrales2023costos}. Their national assessment considers system-level effects, whereas the Bogotá application estimates how household savings vary across tariff strata and simulated price and demand paths. The reported intervals for the change in total cost of ownership remain below zero in all six strata, although the magnitude of the savings varies. One mechanism is the higher efficiency of induction cooking, which reduces the final energy required to deliver the same cooking service, consistent with \citet{im2020} and \citet{scott2024energy}.

Estimated annual emissions reductions are smaller than the normalized value of approximately 0.240~tCO$_2$e per household per year derived from \citet{cabrales2023costos}. In Bogotá, reductions range from 0.064 to 0.151~tCO$_2$e per household per year in the deterministic evaluation and from 0.041 to 0.109~tCO$_2$e across the stochastic means. Differences in baseline consumption, appliance assumptions, and system boundaries may account for part of the gap. This comparison is consistent with the importance that prior studies assign to appliance efficiency, electricity supply, and household conditions \citep{scott2024energy,im2020,floess2023}. In the Bogotá model, lower ownership costs and positive avoided emissions yield negative marginal abatement costs across all strata. The deterministic, mean, and median abatement cost curves also retain the same stratum ordering.

The fiscal results illustrate how cross-subsidy rules can produce different outcomes for households and for the financing of residential energy services. The national analysis by \citet{cabrales2023costos} shows that additional electricity-related fiscal costs can offset reductions in gas subsidies under some assumptions. In the Bogotá base case, avoided subsidy expenditure in Strata~1 and~2 outweighs the increase in Stratum~3 and the loss of contribution revenue from Strata~5 and~6. The resulting expected net fiscal saving is USD~538.26 million in present value over ten years. This balance depends on Bogotá's subscriber composition, consumption, and tariff structure. The findings complement work on the distributional effects of Colombian gas subsidies \citep{cabrales2025subsidios_gas}.

The distribution of household savings also matters for program design. Higher tariff strata show larger absolute lifetime savings, while households in lower strata may face tighter constraints on paying for equipment and installation upfront. The model estimates ownership costs conditional on replacement. Evaluating whether appliance rebates, concessional credit, or on-bill financing would increase adoption requires evidence on household liquidity, installation needs, participation, and program costs.

For regulators and utilities, the fiscal results suggest tracking subsidy obligations and contribution revenues separately in both services as cooking demand shifts. Any proposal to use avoided gas subsidy expenditure to support appliance purchases would need to account for when those savings materialize, how contribution revenues change, and whether the relevant budget arrangements allow those resources to be redirected. These conditions determine how much of the modeled fiscal saving could support an electrification program.

The import scenarios show how higher gas tariffs could change the comparison, given the supply risks discussed by \citet{cabrales2025fedesarrollo_gas_imports}. Relative to the case with no imported gas, the 50\% and 100\% import scenarios raise the cost of continued gas cooking and increase modeled net fiscal savings. Because the sensitivity analysis holds electricity tariffs and consumption paths fixed, these differences isolate the direct effect of the gas tariff shock. A broader assessment would also account for the potential pass-through of higher gas procurement costs to electricity tariffs through gas-fired generation.

Several assumptions delimit the findings. The model holds appliance efficiencies and emission factors fixed and uses a price-elasticity adjustment to represent changes in energy consumption. Adoption, rebound effects, differences in appliance use, and constraints on electricity networks are outside the modeled scope. Applications to other cities would benefit from local evidence on consumption, tariffs, and subscriber composition. Analyses of specific programs would also need to account for adoption and installation costs, changes in grid emissions, and interactions between gas and electricity prices.

%%%%%%%%%%%%%%%%%%%%%%%%%%%%%%%%%%%%%%%%%%%%%%%%%%%%%%%
\section{Conclusions and Policy Implications}
\label{conclusions}
%%%%%%%%%%%%%%%%%%%%%%%%%%%%%%%%%%%%%%%%%%%%%%%%%%%%%%%

This study evaluates residential cooking electrification under uncertain regulated tariffs. The framework links tariff and demand simulations to household ownership costs, greenhouse gas emissions, marginal abatement costs, and changes in subsidies and contributions. The Bogotá application shows how tariff design shapes the distribution of household savings and the fiscal consequences of shifting cooking demand from gas to electricity.

Under the modeled conditions, induction lowers expected ownership costs and annual emissions across all six socioeconomic strata. Expected ten-year savings range from approximately USD~253 per household in Stratum~1 to USD~2,532 in Stratum~6. The reported marginal abatement cost intervals are below zero in every stratum. The aggregate fiscal balance also improves because avoided gas subsidies outweigh the increase in net electricity subsidies and the loss of contribution revenue. Higher gas tariffs strengthen the household and fiscal case for electrification in the import scenarios evaluated.

Two implications follow for utility policy. First, targeted support for appliance purchases and household electrical adaptations could address upfront costs where households have limited access to financing. Second, regulators and utilities should assess the timing and distribution of changes in both subsidy expenditure and contribution revenue as households switch services. The modeled savings provide a basis for evaluating these options. Their implementation depends on adoption, program costs, and local infrastructure and budget conditions.

%%%%%%%%%%%%%%%%%%%%%%%%%%%%%%%%%%%%%%%%%%%%%%%%%%%%%%%
\section*{CRediT authorship contribution statement}

Julio C. Cruz: Conceptualization, Methodology, Software, Formal analysis,
Data curation, Investigation, Visualization, Writing -- original draft,
Writing -- review and editing.

Sergio Cabrales: Conceptualization, Methodology, Supervision, Validation,
Investigation, Writing -- review and editing.

%%%%%%%%%%%%%%%%%%%%%%%%%%%%%%%%%%%%%%%%%%%%%%%%%%%%%%%
\section*{Declaration of competing interest}

The authors declare that they have no known competing financial interests
or personal relationships that could have influenced the work reported
in this paper.

%%%%%%%%%%%%%%%%%%%%%%%%%%%%%%%%%%%%%%%%%%%%%%%%%%%%%%%
\section*{Funding}

This research did not receive any specific grant from funding agencies
in the public, commercial, or not-for-profit sectors.

%%%%%%%%%%%%%%%%%%%%%%%%%%%%%%%%%%%%%%%%%%%%%%%%%%%%%%%
\section*{Data availability}

The study uses publicly available administrative, regulatory, and
macroeconomic data from the Sistema Único de Información de Servicios
Públicos Domiciliarios (SUI), the Comisión de Regulación de Energía y Gas
(CREG), and official Colombian macroeconomic series. The processed data
and code used to reproduce the results are available at
\url{https://github.com/julcruz94-crypto/Evaluating-Residential-Cooking-Electrification-Policies-under-Tariff-Uncertainty}.

%%%%%%%%%%%%%%%%%%%%%%%%%%%%%%%%%%%%%%%%%%%%%%%%%%%%%%%%%%
%%%%%%%%%%%%%%%%%%%%%%%%%%%%%%%%%%%%%%%%%%%%%%%%%%%%%%%%%%

\section*{Declaration of generative AI and AI-assisted technologies in the manuscript preparation process}

During the preparation of this work, the authors used ChatGPT to support language editing and text corrections. After using this tool, the authors reviewed and edited the content as needed and take full responsibility for the content of the published article.

%%%%%%%%%%%%%%%%%%%%%%%%%%%%%%%%%%%%%%%%%%%%%%%%%%%%%%%%%%
%%%%%%%%%%%%%%%%%%%%%%%%%%%%%%%%%%%%%%%%%%%%%%%%%%%%%%%%%%
%%%%%%%%%%%%%%%%%%%%%%%%%%%%%%%%%%%%%%%%%%%%%%%%%%%%%%%%%%
\appendix
\numberwithin{equation}{section}
\numberwithin{figure}{section}
\numberwithin{table}{section}
\renewcommand{\theequation}{\Alph{section}.\arabic{equation}}
\renewcommand{\thefigure}{\Alph{section}.\arabic{figure}}
\renewcommand{\thetable}{\Alph{section}.\arabic{table}}
\section{Technical details of the stochastic tariff and demand simulation}
\label{appendix:stochastic_simulation}
%%%%%%%%%%%%%%%%%%%%%%%%%%%%%%%%%%%%%%%%%%%%%%%%%%%%%%%%%%

This appendix provides the technical details of the stochastic tariff and demand simulation module summarized in Section~\ref{sec:variable_simulation}. The objective of this module is to generate Monte Carlo trajectories for the tariff drivers that concentrate the main uncertainty in residential energy costs: the natural gas supply component ($G_t$), the electricity generation component ($GM_t$), and inflation ($\pi_t$). These simulated trajectories are later translated into regulated end-user tariffs through the tariff formulas defined in Eqs.~\eqref{eq:gas_tariff_method} and \eqref{eq:elec_tariff_method}. The resulting tariff paths are then used to adjust per-user household energy demand through a price-elasticity representation.

%=====================================================
\subsection{Data preparation and log-return construction}
\label{appendix:prep}
%=====================================================

Let $S_t^{(c)}$ denote the monthly level in USD of variable $c \in \{G,GM,\pi\}$ at month $t$. The first step is to prepare the data for stochastic modeling. The database is filtered to retain observations corresponding to the empirical application and to each variable under analysis. The series for $G_t$ and $GM_t$ are expressed in USD using the corresponding exchange rate, and all observations are ordered chronologically. To ensure temporal consistency, the monthly frequency is standardized by retaining the last available observation in each month.

Once the monthly series are consolidated, returns are constructed in logarithmic form. This choice is standard in stochastic price modeling because it is consistent with diffusion-based specifications such as geometric Brownian motion \citep{ibrahim2021geometric} and reduced-form electricity-price models such as Lucia--Schwartz \citep{lucia2002electricity}. Monthly returns are defined as:
\begin{equation}
    r_t^{(c)} = \log\!\left(\frac{S_t^{(c)}}{S_{t-1}^{(c)}}\right),
    \qquad c \in \{G,GM,\pi\},
\end{equation}
for $t = 2,\dots,T$. These returns provide the input for the diagnostic stage and the subsequent choice of simulation models.

%=====================================================
\subsection{Diagnostic tests}
\label{appendix:diagnostics}
%=====================================================

After returns are constructed, their statistical properties are evaluated. This diagnostic block is common to $G_t$, $GM_t$, and $\pi_t$, and determines the simulation route.

Four properties are evaluated. First, serial correlation is tested using the Ljung--Box statistic \citep{ljung1978measure}. For a maximum lag $L$, the test statistic is:
\begin{equation}
    Q_L = T(T+2)\sum_{k=1}^L \frac{\hat \rho_k^2}{T-k},
\end{equation}
where $\hat \rho_k$ is the sample autocorrelation at lag $k$.

Second, conditional heteroscedasticity is evaluated using the ARCH LM test proposed by \citet{engle1982autoregressive}. Third, departures from Normality are assessed through the Jarque--Bera test \citep{jarque1987test} and the Shapiro--Wilk test \citep{shapiro1965analysis}. Fourth, for the pair $(G_t,GM_t)$, the correlation between returns is evaluated in order to determine whether both variables should be simulated jointly or separately.

The use of autocorrelation tests follows the Box--Jenkins identification framework \citep{box2015time}, while the treatment of heteroscedasticity and conditional volatility follows the ARCH/GARCH tradition introduced by \citet{engle1982autoregressive} and extended by \citet{bollerslev1986generalized}. In the case of inflation, the diagnostic results justify a mean-reversion specification. In the case of $G_t$ and $GM_t$, the results determine whether the simulation proceeds through bivariate or univariate models and whether Gaussian or heavy-tailed innovations are more appropriate.

%=====================================================
\subsection{Innovation distribution under non-Normal returns}
\label{appendix:innovation_dist}
%=====================================================

When the Jarque--Bera \citep{jarque1987test} and/or Shapiro--Wilk \citep{shapiro1965analysis} tests reject Normality, Gaussian innovations may understate tail risk and produce predictive intervals that are too narrow. For that reason, the simulation framework allows for Student-\emph{t} innovations in models with conditional heteroscedasticity.

Specifically, when Normality is rejected and ARCH effects are present, innovations are modeled as:
\begin{equation}
z_t \sim t_\nu,
\end{equation}
where $\nu$ denotes the degrees of freedom. If Normality is not rejected, the benchmark specification is maintained:
\begin{equation}
z_t \sim \mathcal{N}(0,1).
\end{equation}

This rule ensures that the model selection process captures potential heavy tails in energy-price returns while preserving a parsimonious specification when Gaussian innovations are adequate.

%=====================================================
\subsection{Inflation modeling}
\label{appendix:inflation}
%=====================================================

Inflation ($\pi_t$) is simulated separately from the two energy components. This treatment is consistent with its role as a macroeconomic driver of tariff indexation rather than as a market-traded energy price.

Inflation is represented through an AR(1) mean-reverting process:
\begin{equation}
\pi_t - \bar{\pi} = \varphi\,(\pi_{t-1} - \bar{\pi}) + \varepsilon_t,
\qquad \varepsilon_t \sim \mathcal{N}(0,\sigma_\pi^2),
\end{equation}
where $\bar{\pi}$ is the long-run mean, $\varphi$ measures persistence, and $|\varphi|<1$ guarantees mean reversion. The parameters $(\bar{\pi},\varphi,\sigma_\pi)$ are estimated using the historical monthly inflation series. Future inflation paths are generated recursively over the simulation horizon and later used in the tariff reconstruction stage. This specification is consistent with standard representations of inflation dynamics in macroeconomic time series, where persistence and mean reversion are commonly captured through low-order autoregressive processes, as in \citet{khamis2020} and \citet{babu2015}.

%=====================================================
\subsection{Correlation-based simulation strategy for $G_t$ and $GM_t$}
\label{appendix:bivariate}
%=====================================================

The simulation of the natural gas supply component ($G_t$) and the electricity generation component ($GM_t$) depends on the empirical relationship between their return series. Let $r_t^{(G)}$ and $r_t^{(GM)}$ denote the synchronized monthly returns of the gas and electricity components. Their empirical dependence is evaluated using the Pearson correlation coefficient:
\begin{equation}
    \rho_{G,GM} = \mathrm{corr}\bigl(r_t^{(G)}, r_t^{(GM)}\bigr), 
    \qquad H_0: \rho_{G,GM} = 0.
\end{equation}

If the null hypothesis is rejected, $G_t$ and $GM_t$ are simulated jointly. If the null cannot be rejected, they are simulated separately. In both cases, the Normality diagnostic determines the admissible innovation structure.

This approach is consistent with the joint modeling of energy prices used to capture co-movements and substitution effects in studies such as \citet{tan2010day} and \citet{ibrahim2021geometric}.

%=====================================================
\subsection{Correlated case: joint simulation of $G_t$ and $GM_t$}
\label{appendix:dependent_cases}
%=====================================================

If gas and electricity returns are significantly correlated, the model proceeds with joint simulation of $G_t$ and $GM_t$. At this point, the Normality diagnostic determines the candidate models.

If Normality is not rejected, the simulation is restricted to bivariate models, and the preferred specification is selected using AIC and BIC. Two candidate models are considered.

The first is a bivariate geometric Brownian motion:
\begin{equation}
        \begin{pmatrix} r_h^{(G)} \\ r_h^{(GM)} \end{pmatrix}
        =
        \begin{pmatrix} \mu_G \\ \mu_{GM} \end{pmatrix}
        +
        \begin{pmatrix} \sigma_G & 0 \\ 0 & \sigma_{GM} \end{pmatrix}
        L \mathbf{z}_h,\qquad \mathbf{z}_h \sim \mathcal{N}(\mathbf{0}, I_2),
\end{equation}
where $L$ is the Cholesky factor of the correlation matrix and the means, volatilities, and correlation are estimated from the data. This type of multivariate lognormal representation is commonly used in energy and financial modeling to capture correlated stochastic dynamics between related commodities \citep{ibrahim2021geometric}.

The second is a regression-based dependence model:
\begin{equation}
        r_t^{(GM)} = \alpha + \beta\, r_t^{(G)} + e_t,
        \qquad e_t \sim \mathcal{N}(0,\sigma_e^2),
\end{equation}
where gas returns are simulated first and electricity returns are generated conditionally. This specification captures conditional dependence between energy prices and is consistent with reduced-form approaches used in energy price modeling such as those discussed by \citet{tan2010day}.

If returns are correlated but Normality is rejected, the full model set is evaluated. In this branch, model choice is again based on AIC and BIC, but the admissible set is broader and includes specifications that allow heavy tails, autoregressive dynamics, and seasonality. The candidate models for the electricity component are:

\begin{enumerate}
    \item A regression-based dependence model:
    \begin{equation}
        r_t^{(GM)} = \alpha + \beta\, r_t^{(G)} + e_t,
    \end{equation}
    where $e_t$ may follow a Student-\emph{t} distribution.

    \item A univariate AR--GARCH model:
    \begin{align}
        r_t^{(GM)} &= \alpha_0 + \sum_{j=1}^p \phi_j r_{t-j}^{(GM)} + \varepsilon_t,\\
        \varepsilon_t &= \sigma_t z_t,\qquad z_t \sim \mathcal{D}_{GM},\\
        \sigma_t^2 &= \omega + \alpha \varepsilon_{t-1}^2 + \beta \sigma_{t-1}^2.
    \end{align}

    \item A Fourier-based seasonal model:
    \begin{equation}
        r_t^{(GM)} = f_{\text{Fourier}}(m_t) + \varepsilon_t,
    \end{equation}
    where $m_t$ denotes the calendar month.

    \item A monthly dummy regression:
    \begin{equation}
        r_t^{(GM)} = \alpha_0 + \sum_{m=2}^{12} \gamma_m \mathbb{I}\{m_t=m\} + \varepsilon_t.
    \end{equation}

    \item A Lucia--Schwartz-type reduced-form model inspired by \citet{lucia2002electricity}:
    \begin{equation}
        \log S_t^{(GM)} = f_t + x_t,
    \end{equation}
    with
    \begin{equation}
        x_t = a + b x_{t-1} + \varepsilon_t,
        \qquad \varepsilon_t \sim \mathcal{N}(0,\sigma_\varepsilon^2),
    \end{equation}
    where $|b|<1$.
\end{enumerate}

Seasonal representations based on Fourier terms and calendar dummies are commonly used to capture periodic patterns in energy markets and electricity prices, as in \citet{lucia2002electricity}. More generally, the AR--GARCH specification follows the standard conditional-volatility framework introduced by \citet{engle1982autoregressive} and generalized by \citet{bollerslev1986generalized}, while the autoregressive identification logic is consistent with the Box--Jenkins approach described by \citet{box2015time}.

All candidate models are estimated on the training sample, ranked using AIC and BIC, and the preferred specification is then re-estimated on the full sample for simulation.

%=====================================================
\subsection{Train--test split and information criteria}
\label{appendix:model_selection_dep}
%=====================================================

When more than one specification is admissible, model selection is based on an out-of-sample train--test split. Approximately 80\% of the observations are used for estimation and 20\% are reserved for validation.

For each model, the test-sample log-likelihood $\ell_{\text{test}}$ is computed, and the information criteria are defined as:
\begin{align}
    \mathrm{AIC} &= 2k - 2\ell_{\text{test}},\\
    \mathrm{BIC} &= k \log n - 2\ell_{\text{test}},
\end{align}
where $k$ is the number of parameters and $n$ is the number of observations in the test sample. These criteria follow the standard formulations proposed by \citet{akaike2003new} and \citet{schwarz1978estimating}. The preferred specification is the one with the lowest BIC.

After the preferred model is selected, it is subjected to an out-of-sample backtest. The parameters estimated on the training sample are used to simulate a distribution of paths over the validation window. The observed trajectory of $GM_t$ is then compared with the predictive intervals implied by the model. This provides an empirical check that the selected specification reproduces the observed dynamics with reasonable accuracy.

Additional graphical diagnostics are also produced, including correlation matrices in levels and returns, and scatter plots of $(G_t,GM_t)$ and $(r_t^{(G)},r_t^{(GM)})$. These graphs are used for interpretation, but the formal selection rule remains based on BIC.

%=====================================================
\subsection{Independent simulation of $G_t$ and $GM_t$}
\label{appendix:independent_simulation}
%=====================================================

If the correlation test does not reject $H_0:\rho_{G,GM}=0$, $G_t$ and $GM_t$ are simulated separately. In this branch, the Normality diagnostic determines whether the simulation uses standard univariate models or univariate models with heavy-tailed volatility.

The simulation horizon is set to $H$ months, with $H=120$, and $N$ Monte Carlo paths are generated for each component.

%-----------------------------------------------------
\subsubsection{Independent case with approximately Normal returns}
\label{appendix:independent_normal}
%-----------------------------------------------------

If returns are approximately Normal, the simulation proceeds with standard univariate models. Depending on the diagnostics for each series, these may take the form of geometric Brownian motion or autoregressive dynamics.

For a component $c \in \{G,GM\}$, if the process is approximated by an independent return specification, future returns are generated as:
\begin{equation}
    r_{t+h}^{(c)} = \mu_c + \sigma_c Z_h,
    \qquad Z_h \sim \mathcal{N}(0,1),
\end{equation}
and the corresponding price paths are reconstructed as:
\begin{equation}
    S_{t+h}^{(c)} = S_t^{(c)} \exp\!\left( \sum_{j=1}^h r_{t+j}^{(c)} \right).
\end{equation}

If serial dependence is present, the return process is represented through an autoregression:
\begin{equation}
    r_t^{(c)} = \alpha_0 + \sum_{j=1}^p \phi_j r_{t-j}^{(c)} + \varepsilon_t,
    \qquad \varepsilon_t \sim \mathcal{N}(0,\sigma_\varepsilon^2).
\end{equation}
Future returns are then generated recursively from the estimated process. These specifications are consistent with the autoregressive representation of time-series dynamics developed in the Box--Jenkins framework \citep{box2015time}.

%-----------------------------------------------------
\subsubsection{Independent case with non-Normal returns}
\label{appendix:independent_nonnormal}
%-----------------------------------------------------

If returns are not Normal, the simulation continues through univariate models with GARCH-type volatility and Student-\emph{t} innovations. In this branch, the basic return equation is:
\begin{align}
    r_t^{(c)} &= \mu_c + \varepsilon_t, \\
    \varepsilon_t &= \sigma_t z_t,\qquad z_t \sim \mathcal{D}_{c}, \\
    \sigma_t^2 &= \omega + \alpha \varepsilon_{t-1}^2 + \beta \sigma_{t-1}^2,
\end{align}
where $\mathcal{D}_{c}=t_\nu$ when heavy tails are required.

If serial dependence is also present, the autoregressive mean is combined with GARCH volatility:
\begin{align}
    r_t^{(c)} &= \alpha_0 + \sum_{j=1}^p \phi_j r_{t-j}^{(c)} + \varepsilon_t,\\
    \varepsilon_t &= \sigma_t z_t,\qquad z_t \sim \mathcal{D}_{c},\\
    \sigma_t^2 &= \omega + \alpha \varepsilon_{t-1}^2 + \beta \sigma_{t-1}^2.
\end{align}

Future paths are simulated recursively by propagating both the conditional mean and the conditional variance. In all independent cases, the result is a collection of simulated trajectories $\{S_{t+h}^{(c,i)}\}$ for $i=1,\dots,N$. These specifications follow the ARCH framework proposed by \citet{engle1982autoregressive}, its GARCH extension by \citet{bollerslev1986generalized}, and the broader autoregressive logic described by \citet{box2015time}.

%=====================================================
\subsection{Tariff reconstruction under the regulatory framework}
\label{appendix:tariff_reconstruction}
%=====================================================

Once the stochastic paths for $G_t$, $GM_t$, and $\pi_t$ are generated, they are translated into regulated end-user tariffs. This step follows the tariff formulas defined by the regulatory framework corresponding to the empirical application.

Let $CU_t$ denote the residential electricity tariff in USD/kWh and $Cuv_t$ the residential natural gas tariff in USD/m$^3$. For each Monte Carlo path $i$ and month $t$, the simulated values $\{G_t^{(i)}, GM_t^{(i)}, \pi_t^{(i)}\}$ are introduced into the regulated tariff equations. Stochasticity therefore enters through these three drivers, while the remaining tariff components evolve deterministically according to regulatory rules, indexation mechanisms, and caps.

This step produces a set of simulated end-user tariffs $\{CU_t^{(i)}, Cuv_t^{(i)}\}$, which are used in the cost calculations developed in the main methodology.

%=====================================================
\subsection{Outputs for downstream techno-economic analysis}
\label{appendix:outputs}
%=====================================================

The output of the stochastic module is a set of Monte Carlo trajectories for monthly electricity and gas tariffs in USD. These trajectories summarize the range of possible future price paths under the regulatory framework defined in the empirical application. In the subsequent stages of the methodology, they are used to estimate stochastic total cost of ownership, greenhouse gas emissions, and marginal abatement costs for residential cooking electrification.

This use of Monte Carlo simulation is consistent with standard approaches for propagating uncertainty in stochastic frameworks applied to energy systems and economic evaluation, such as those used by \citet{khamis2020} and \citet{babu2015}.

%=====================================================
\subsection{Price-responsive household demand adjustment}
\label{appendix:demand_adjustment}
%=====================================================

After the simulated tariff paths are reconstructed, they are used to adjust household energy demand through a price-elasticity representation. This step introduces stochastic variation in per-user demand and allows emissions and abatement potential to vary across Monte Carlo paths.

The demand adjustment follows the standard log-linear specification used in residential energy demand studies:
\begin{equation}
\ln Q_{x,t}^{(i)} = \alpha_x + \varepsilon_x \ln P_{x,t}^{(i)} + u_{x,t}^{(i)},
\label{eq:appendix_loglinear_demand}
\end{equation}
where $Q_{x,t}^{(i)}$ is the household consumption of energy carrier $x$ in month $t$ and simulation path $i$, $P_{x,t}^{(i)}$ is the corresponding simulated end-user price, $\varepsilon_x$ is the price elasticity of residential energy demand, $\alpha_x$ is the intercept, and $u_{x,t}^{(i)}$ is the error term.

In this paper, the elasticity parameter is not estimated from household-level panel data. Instead, it is taken from the residential electricity demand literature. In particular, \citet{espey2004turning} provide a meta-analysis of residential electricity demand elasticities and show that residential demand responds negatively to price changes, with stronger responses in the long run than in the short run. This evidence supports the use of an elasticity-based adjustment when the objective is to propagate price-induced demand variation into cost and emissions outcomes without estimating a complete behavioral demand system.

For implementation, Eq.~\eqref{eq:appendix_loglinear_demand} is expressed relative to the baseline household consumption and baseline tariff:
\begin{equation}
Q_{x,t}^{(i)} =
Q_{x,0}
\left(
\frac{P_{x,t}^{(i)}}{P_{x,0}}
\right)^{\varepsilon_x},
\label{eq:elasticity_demand}
\end{equation}
where $Q_{x,0}$ is the baseline household consumption of energy carrier $x$, and $P_{x,0}$ is the corresponding baseline tariff.

Because residential energy demand is price-inelastic, $\varepsilon_x<0$. Therefore, when the simulated tariff $P_{x,t}^{(i)}$ increases relative to the baseline tariff $P_{x,0}$, the adjusted demand $Q_{x,t}^{(i)}$ decreases. Conversely, when the simulated tariff decreases, the adjusted demand increases. This mechanism allows the model to reflect the direction of household demand response to price changes while maintaining a parsimonious structure.

The adjusted demand is then used in the operating-cost calculation:
\begin{equation}
OPEX_{x,t}^{(i)} = Price_{x,t}^{(i)} \times Consumption_{x,t}^{(i)},
\end{equation}
and in the emissions calculation:
\begin{equation}
GHG_x^{(i)} =
EF_x \times StoveConsumption_x^{(i)}.
\end{equation}

Thus, tariff uncertainty affects the results through two channels. First, it directly changes the price paid by households. Second, it indirectly modifies household consumption through the elasticity-based demand adjustment. This makes TCO, GHG emissions, abatement potential, and MAC values stochastic outcomes.

%%%%%%%%%%%%%%%%%%%%%%%%%%%%%%%%%%%%%%%%%%%%%%%%%%%%%%%%%%
%%%%%%%%%%%%%%%%%%%%%%%%%%%%%%%%%%%%%%%%%%%%%%%%%%%%%%%%%%
\section{Additional case-study context}
\label{appendix:case_context}
%%%%%%%%%%%%%%%%%%%%%%%%%%%%%%%%%%%%%%%%%%%%%%%%%%%%%%%%%%

This appendix reports additional descriptive evidence used to characterize the Colombian residential natural gas and electricity context. These figures support the case-study discussion but are not required for the main interpretation of the TCO, fiscal, emissions, and MAC results.

%=====================================================
\subsection{National natural gas context}
\label{appendix:national_gas_context}
%=====================================================

The residential natural gas market analyzed in this paper is part of a broader Colombian gas system that includes domestic production, imports, transportation, distribution, and final consumption across multiple sectors. Figure~\ref{fig:natgas-donut} shows the average distribution of natural gas consumption by sector for 2006--2024. The industrial sector accounts for the largest share of demand, while the residential sector represents 26.9\% of total consumption. This share is relevant because residential gas use is not a marginal component of the national gas system. Therefore, policies that affect residential cooking technologies can have implications beyond individual households, including effects on gas demand, tariff revenues, subsidy flows, and exposure to gas-supply risks.

\begin{figure}[H]
\centering
\includegraphics[width=0.75\linewidth]{NatGasCons.png}
\caption{\textbf{Natural gas consumption by sector in Colombia (average 2006--2024).}
The figure presents the sectoral distribution of national natural gas demand, highlighting the role of the industrial and residential sectors, followed by transport, commercial and public consumption, mining, and other uses. Source: authors' elaboration based on \citet{UPME_BECO_2026}.}
\label{fig:natgas-donut}
\end{figure}

Figure~\ref{fig:supply_chain} summarizes the structure of the Colombian natural gas supply chain. Domestic production and imports enter the system upstream and converge before final distribution to residential, commercial, industrial, transport, refinery, thermal power, and compressor-station users. This structure matters for the proposed framework because residential tariffs are not treated as purely exogenous market prices. Instead, they are interpreted as outcomes of regulated formulas that transfer supply, transportation, distribution, commercialization, technical adjustments, and macroeconomic indexation to end users.

\begin{figure}[H]
\centering
\includegraphics[width=\linewidth]{\detokenize{Gas Supply Chain.png}}
\caption{\textbf{Natural gas supply chain in Colombia.}
Domestic production and imports converge before the citygate and distribution stage, from which gas is delivered to residential, commercial, transport, refinery, thermal power, industrial, and compressor-station users. Source: authors' elaboration based on \citet{cabrales2023costos}.}
\label{fig:supply_chain}
\end{figure}

The upstream structure shown in Figure~\ref{fig:supply_chain} is important for interpretation because changes in gas supply conditions, including increased reliance on imports, may be transmitted to households through the regulated tariff structure.

%=====================================================
\subsection{Residential gas subscribers}
\label{appendix:residential_gas_subscribers}
%=====================================================

Administrative records on residential gas subscribers are used to characterize both the scale and the spatial distribution of residential gas use in Colombia. Subscriber counts vary significantly across territories, reflecting differences in infrastructure deployment, demographic concentration, and urban development.

Over the last two decades, Colombia's residential natural gas network has expanded substantially \citep{cabrales2025subsidios_gas}. At the same time, the distribution of connected households remains strongly concentrated in a limited set of departments where urban density and gas infrastructure coincide. Figure~\ref{fig:gas-users-department} illustrates this pattern by comparing the number of residential gas users by department in 2005 and in the last year with available data.

\begin{figure}[H]
\centering
\includegraphics[width=0.90\linewidth]{HeatMapGasUsers.png}
\caption{\textbf{Residential gas connections by department in Colombia (2005 vs.\ last year with available data).}
The maps report the total number of connected residential gas users aggregated at the departmental level. Darker shades indicate a higher number of users. The blue marker highlights the location of Bogotá, which serves as the benchmark case study. Source: authors' elaboration based on SUI data.}
\label{fig:gas-users-department}
\end{figure}

Figure~\ref{fig:gas-users-department} shows two stylized facts. First, residential gas service expanded geographically between 2005 and the latest year for which data are available, reflecting the progressive extension of the distribution network across the country. Second, despite this expansion, the bulk of residential subscribers remains concentrated around the largest urban systems and their surrounding municipalities. In practice, regulated residential gas demand is concentrated in metropolitan areas such as Barranquilla, Bogotá, Bucaramanga, Cali, Cartagena, Cúcuta, Ibagué, Medellín, Soacha, and Soledad.

In addition to its geographic concentration, the residential gas market shows a clear pattern of sustained growth together with marked socioeconomic stratification, as shown in Figures~\ref{fig:national-users} and \ref{fig:donut-strata}. The number of regulated residential users increased from approximately 3.57 million in 2005 to more than 10.41 million in 2022 (Figure~\ref{fig:national-users}). This expansion is largely explained by the progressive rollout of gas infrastructure and by policies aimed at increasing access, particularly among lower-income households. Unless otherwise stated, the descriptive figures reported in this appendix are based on official SUI data processed by the authors and aligned to the December~2022 snapshot.

\begin{figure}[H]
\centering
\includegraphics[width=0.98\linewidth]{AnnualResGas.png}
\caption{Residential gas subscribers (national, December of each year).}
\label{fig:national-users}
\end{figure}

\begin{figure}[H]
\centering
\includegraphics[width=0.64\linewidth]{ResGasSuscribers.jpg}
\caption{Residential gas subscribers by stratum (snapshot, 2022--12).}
\label{fig:donut-strata}
\end{figure}

Two elements are particularly important for the empirical strategy. First, the steady expansion of the residential user base. Second, the tariff structure based on socioeconomic stratification. This setup makes stratum~4 a natural benchmark, as it allows the identification of regulated cost signals without the direct effect of subsidies or surcharges. From there, the analysis can be extended to assess distributional effects across subsidized and contributing strata.

%=====================================================
\subsection{Additional energy-cost comparisons}
\label{appendix:additional_energy_costs}
%=====================================================

Residential gas prices differ across cities because supply conditions, transportation costs, and local distribution costs vary. To facilitate cross-energy comparisons, the residential gas unit cost (Cuv) is expressed on a common energy basis in USD/MMBtu.

Figure~\ref{fig:cons-price-city} jointly presents consumption per subscriber and the average residential gas unit cost, Cuv, in USD/MMBtu for the ten cities with the largest residential gas-user base. This joint representation highlights that higher consumption does not necessarily coincide with higher unit costs. In the results section, the baseline estimates are reported first for Bogotá stratum~4 and then extended through comparisons across strata and cities.

\begin{figure}[H]
\centering
\includegraphics[width=0.92\linewidth]{GasConsumptionandPrice.png}
\caption{Gas consumption per subscriber (bars) and average unit cost—Cuv (dots) by city (2022--12) for the ten cities with the largest residential gas user base. Unit cost is expressed in USD/MMBtu.}
\label{fig:cons-price-city}
\end{figure}

A direct comparison between natural gas and electricity tariffs is not possible in their original physical units, since one is charged in USD/m$^3$ and the other in USD/kWh. For that reason, both series are converted to a common energy basis in USD/MMBtu. Figure~\ref{fig:energy-cost-comparison} reports the residential gas tariff (Cuv) and the residential electricity tariff (CU) in USD/MMBtu for Bogotá stratum~4 without efficiency adjustment. This normalization allows a direct comparison of purchased energy costs over time.

\begin{figure}[H]
\centering
\includegraphics[width=0.95\linewidth]{ComparaGasElecSinCorr.png}
\caption{\textbf{Comparison of residential energy costs in Bogotá (USD/MMBtu).}
Monthly evolution of natural gas (Cuv) and electricity (CU) tariffs converted to a common energy basis.}
\label{fig:energy-cost-comparison}
\end{figure}

Table~\ref{tab:energy-cost-bogota} complements the graphical evidence with descriptive statistics for the Bogotá stratum~4 benchmark. Both natural gas and electricity exhibit time variation, but electricity is systematically more expensive on a common purchased-energy basis. This persistent gap is one of the central empirical features that motivates the evaluation of electrification through both efficiency corrections and stochastic tariff simulation.

\begin{table}[H]
\centering
\caption{\textbf{Descriptive statistics of residential energy unit costs in Bogotá (stratum~4)} expressed in USD/MMBtu (non-efficiency-adjusted basis).}
\label{tab:energy-cost-bogota}
\resizebox{0.98\linewidth}{!}{
\begin{tabular}{lrrrrrrrrr}
\toprule
\textbf{Energy carrier} & \textbf{Count} & \textbf{Mean} & \textbf{Min} & \textbf{Max} & \textbf{Std} & \textbf{Q1} & \textbf{Q3} & \textbf{Skew.} & \textbf{Kurt.} \\
\midrule
Natural gas (Cuv) & 125 & 16.02 & 12.22 & 23.79 & 2.50 & 14.51 & 16.37 & 1.4445 & 1.9369 \\
Electricity (CU)  & 150 & 41.27 & 27.60 & 68.56 & 9.92 & 32.88 & 46.54 & 0.9354 & 0.1806 \\
\bottomrule
\end{tabular}}
\end{table}

%=====================================================
\subsection{Additional regulated tariff details}
\label{appendix:additional_tariff_details}
%=====================================================

Residential natural gas tariffs in Colombia are defined and periodically updated according to CREG methodologies. Transportation charges ($T_m$) are typically indexed to the Producer Price Index (IPP), distribution costs ($D_m$) and commercialization-related components may be updated through Consumer Price Index (IPC) indexation, and the supply component ($G_m$) reflects procurement conditions that may be exposed to international price dynamics and exchange-rate movements. Technical parameters such as the loss factor ($\rho_m$) and calorific adjustment ($fpc_m$) are also embedded in the tariff formulation \citep{CREG137_2013,CREG183_2013,CREG184_2014}.

Figure~\ref{fig:price-gas-bog-e4} reports the monthly evolution of the residential natural gas unit cost (Cuv) in Bogotá for stratum~4, expressed in USD/MMBtu. Because stratum~4 is not affected by subsidies or contributions, this series approximates the underlying regulated gas-cost signal used in the baseline analysis.

\begin{figure}[H]
\centering
\includegraphics[width=0.95\linewidth]{PGasBogota.png}
\caption{\textbf{Monthly evolution of the residential natural gas unit cost (Cuv) in Bogotá (stratum~4).}
The figure reports the historical trajectory of Cuv expressed in USD/MMBtu for the benchmark case study used in the baseline analysis.}
\label{fig:price-gas-bog-e4}
\end{figure}

Residential electricity tariffs are defined through a multi-component structure that includes generation, transmission, distribution, commercialization, recognized losses, and restriction costs \citep{CREG119_2007}. The generation component ($GM_m$) reflects energy procurement through wholesale markets and bilateral contracts, while network and commercialization components are regulated and may incorporate macroeconomic indexation. These mechanisms imply that electricity tariffs evolve through both market-based and regulated channels.

Figure~\ref{fig:price-elec-bog-e4} reports the monthly evolution of the residential electricity unit cost (CU) in Bogotá for stratum~4, expressed in USD/MMBtu. As in the gas case, focusing on stratum~4 helps isolate the regulated cost signal from the direct effects of subsidies and contributions.

\begin{figure}[H]
\centering
\includegraphics[width=0.95\linewidth]{PElecBog.png}
\caption{\textbf{Monthly evolution of the residential electricity unit cost (CU) in Bogotá (stratum~4).}
The figure reports the historical trajectory of CU expressed in USD/MMBtu for the benchmark case study used in the baseline analysis.}
\label{fig:price-elec-bog-e4}
\end{figure}

%%%%%%%%%%%%%%%%%%%%%%%%%%%%%%%%%%%%%%%%%%%%%%%%%%%%%%%%%%
\section{Stochastic diagnostics and model-selection results}
\label{appendix:stochastic_diagnostics}
%%%%%%%%%%%%%%%%%%%%%%%%%%%%%%%%%%%%%%%%%%%%%%%%%%%%%%%%%%

Table~\ref{tab:sim_corr_diag} summarizes the main diagnostic results. Two findings are especially relevant. First, the log-returns of $G_t$ and $GM_t$ show statistically significant dependence, which supports the use of a correlated specification rather than treating both processes as independent. Second, the adequacy tests indicate that the gas component shows no strong evidence of autocorrelation or conditional heteroscedasticity, but does depart from Normality. This combination supports representing $G_t$ as an approximately independent process under a relaxed GBM-type structure with heavy-tailed innovations. 

For the electricity generation component, the diagnostics reported in Table~\ref{tab:sim_corr_diag} lead to a different implication. Although the series also satisfies the relaxed adequacy conditions for autocorrelation and heteroscedasticity, the rejection of Normality and the observed dependence with $G_t$ suggest that $GM_t$ should not be modeled as a simple standalone process. Instead, the results justify evaluating a broader set of dependent and reduced-form alternatives, as established in the methodology.

\begin{table}[H]
\centering
\caption{Dependence and diagnostic results for the stochastic tariff drivers in Bogotá}
\label{tab:sim_corr_diag}
\begin{tabular}{llc}
\toprule
\textbf{Block} & \textbf{Indicator} & \textbf{Result} \\
\midrule
\multirow{1}{*}{Dependence between $G_t$ and $GM_t$} 
& Pearson correlation of log-returns & 0.374 \\
\midrule
\multirow{6}{*}{Natural gas supply component ($G_t$)} 
& Monthly volatility, $\sigma_{monthly}$ & 0.017543 \\
& Annualized reference volatility, $\sigma_{annualized}$ & 0.060770 \\
& Ljung--Box $p$-value & 0.3960 \\
& ARCH test $p$-value & 0.2558 \\
& Jarque--Bera $p$-value & 0.0076 \\
& Shapiro--Wilk $p$-value & 0.0000 \\
\midrule
\multirow{4}{*}{Electricity generation component ($GM_t$)} 
& Ljung--Box $p$-value & 0.9150 \\
& ARCH test $p$-value & 0.1096 \\
& Jarque--Bera $p$-value & 0.0040 \\
& Shapiro--Wilk $p$-value & 0.0000 \\
\bottomrule
\end{tabular}
\end{table}

Table~\ref{tab:model_selection_gm} reports the out-of-sample comparison of the candidate models evaluated for $GM_t$. The ranking shows that the regression-based dependence model with gas prices (REG\_GBM) dominates the alternative specifications according to the information-criterion comparison, and particularly under the BIC decision rule adopted in this paper. This result indicates that, in the Bogotá case, the electricity generation component is better represented when its dynamics are jointly modeled with the natural gas supply component than when it is approximated purely seasonal or standalone reduced-form alternatives.

In substantive terms, the selected model implies that the co-movement between both tariff drivers contains relevant predictive information. This is consistent with the empirical evidence reported in Table~\ref{tab:sim_corr_diag} and provides the basis for using a dependence-based specification in the stochastic simulation stage.

\begin{table}[H]
\centering
\caption{Model comparison for the electricity generation component ($GM_t$) in Bogotá}
\label{tab:model_selection_gm}
\begin{tabular}{lccc}
\toprule
\textbf{Candidate model} & \textbf{Log-likelihood} & \textbf{AIC} & \textbf{BIC} \\
\midrule
REG\_GBM        & 80.63 & -155.25 & -149.77 \\
FOURIER\_UNIV   & 78.34 & -148.67 & -141.36 \\
DUMMIES\_UNIV   & 77.06 & -128.13 & -104.36 \\
LS\_OU\_1F      & 63.78 & -121.56 & -116.07 \\
\bottomrule
\end{tabular}
\end{table}

Figure~\ref{fig:backtest-gm} provides a graphical validation of the selected specification for $GM_t$. The observed trajectory in the test window falls within the simulated prediction interval, supporting the adequacy of the chosen model in out-of-sample terms. Beyond reproducing the series' central tendency, the selected specification also captures the range of variation observed during the validation period, including episodes of higher values near the end of the sample. This result reinforces the model-selection evidence reported in Table~\ref{tab:model_selection_gm}.

\begin{figure}[H]
    \centering
    \includegraphics[width=0.78\textwidth]{BacktestSimGenElec.png}
    \caption{\textbf{Graphical backtest of the selected simulation model for the electricity generation component ($GM_t$).} The observed values in the test period fall within the simulated prediction intervals, supporting the adequacy of the model selected based on AIC/BIC criteria.}
    \label{fig:backtest-gm}
\end{figure}

Figure~\ref{fig:sim-g-gm} presents the historical series and forward simulations for both tariff drivers. In the natural gas case, the simulated paths preserve the broad scale historical variation while showing a moderate widening of the uncertainty band over time. This behavior is consistent with the earlier diagnostics: the process shows little persistence or conditional volatility, but it does require a non-Gaussian innovation structure.

For the electricity generation component, the simulated band is visibly wider. This reflects the larger uncertainty associated with electricity generation costs in Bogotá and the fact that the preferred specification explicitly incorporates the dependence between $GM_t$ and $G_t$. From an economic standpoint, this result is especially important because electricity prices determine the main operating-cost component of the electrification scenario. Hence, the wider uncertainty band for $GM_t$ is one of the main channels through which stochastic variation propagates into TCO and MAC outcomes.

\begin{figure}[H]
    \centering
    \begin{subfigure}[b]{0.49\textwidth}
        \centering
        \includegraphics[width=\textwidth]{SimGenGas.png}
        \caption{Natural gas supply component ($G_t$).}
    \end{subfigure}
    \hfill
    \begin{subfigure}[b]{0.49\textwidth}
        \centering
        \includegraphics[width=\textwidth]{SimGenElec.png}
        \caption{Electricity generation component ($GM_t$).}
    \end{subfigure}
    \caption{\textbf{Historical series and simulated trajectories for the main tariff drivers in Bogotá.}}
    \label{fig:sim-g-gm}
\end{figure}

%%%%%%%%%%%%%%%%%%%%%%%%%%%%%%%%%%%%%%%%%%%%%%%%%%%%%%%%%%
\section{Additional fiscal, MACC, and imported-gas scenario results}
\label{appendix:fiscalMACCandimportedgas}
%%%%%%%%%%%%%%%%%%%%%%%%%%%%%%%%%%%%%%%%%%%%%%%%%%%%%%%%%%

Table~\ref{tab:fiscal_impact} reveals a clear redistribution pattern across socioeconomic strata. Under BAU, lower-income households represent a fiscal cost because they receive subsidies, while higher-income households generate fiscal revenues through contributions. This is reflected in the positive BAU values for Strata 1 and 2, the neutral position of Strata 3 and 4, and the negative BAU values for Strata 5 and 6.

Once electrification is introduced, the fiscal effect changes direction depending on the stratum. For Strata 1 and 2, electrification generates fiscal savings, as indicated by the negative deterministic and stochastic values. This means that part of the subsidy burden currently associated with gas consumption is reduced when households switch to induction cooking. By contrast, for Strata 5 and 6, electrification produces a positive fiscal effect, representing forgone public revenues from the reduction in contribution payments previously associated with gas consumption. Strata 3 and 4 remain close to fiscal neutrality, as expected from their position in the tariff structure.

\begin{table}[H]
\centering
\caption{Fiscal impact per household of cooking electrification under different scenarios}
\label{tab:fiscal_impact}

\resizebox{\textwidth}{!}{
\begin{tabular}{lcccc}
\toprule
\textbf{Stratum} & \textbf{BAU} & \textbf{Deterministic} & \textbf{Stochastic mean} & \textbf{Stochastic (95\% CI)} \\
\textit{Units} & (USD/user) & (USD/user) & (USD/user) & (USD/user) \\
\midrule
E1 & 676.36 & -573.24 & -853.58 & [-964.73,\,-744.54] \\
E2 & 473.66 & -400.17 & -620.95 & [-701.81,\,-541.63] \\
E3 & 0.00 & 18.91 & 22.73 & [17.82,\ 29.39] \\
E4 & 0.00 & 0.00 & 0.00 & [0.00,\ 0.00] \\
E5 & -296.20 & 257.13 & 310.64 & [270.96,\ 351.09] \\
E6 & -950.29 & 890.62 & 474.83 & [389.11,\ 562.46] \\
\bottomrule
\end{tabular}
}
\end{table}

The deterministic curve in Figure~\ref{fig:macc-deter} already shows this result. The absolute magnitude of the MAC increases as one moves from lower to higher strata, reflecting the fact that households with higher baseline consumption generate larger savings per unit of abatement. In the stochastic mean and stochastic median curves, this pattern remains essentially unchanged. The ordering of strata is stable, which means that uncertainty does not alter the relative ranking of the intervention across the socioeconomic distribution.

Another important feature is the cumulative-abatement structure of the MACC. Higher strata contribute more strongly to total abatement potential because their annual emissions reductions are larger. As a result, the right-hand side of the curve is dominated by Stratum 5 and Stratum 6. In policy terms, this means that upper-income strata play a disproportionate role in aggregate emissions reduction, even though electrification remains economically attractive across the entire distribution.

The similarity between Figures~\ref{fig:macc-mean} and \ref{fig:macc-p50} is also informative. Since the mean-based and median-based MACCs are close to each other, the stochastic distributions do not appear to be heavily skewed in a way that would overturn the central interpretation of the results. Instead, uncertainty mainly changes the scale of the bars, while leaving the qualitative shape of the curve intact.

\begin{figure}[H]
    \centering
    \includegraphics[width=0.88\textwidth]{MACCBogDeter.png}
    \caption{\textbf{Deterministic MACC based on historical data.}}
    \label{fig:macc-deter}
\end{figure}
\begin{figure}[H]
    \centering
    \includegraphics[width=0.88\textwidth]{MACCBogP50Elec.png}
    \caption{\textbf{MACC based on median (P50) stochastic MAC values.}}
    \label{fig:macc-p50}
\end{figure}

Table~\ref{tab:import_scenarios_tco_det} shows that higher import dependence substantially increases the cost of remaining in the BAU gas-cooking alternative. Under the 50\% imported-gas scenario, deterministic $\Delta TCO$ ranges from USD $-590.80$ per household in Stratum~1 to USD $-3970.30$ in Stratum~6. Under full import dependence, the savings become even larger, ranging from USD $-894.81$ in Stratum~1 to USD $-5406.25$ in Stratum~6. These results confirm that gas-import exposure strengthens the private economic case for induction cooking across all socioeconomic strata.

\begin{table}[htbp]
\centering
\caption{TCO under imported-gas scenarios in Bogotá: BAU gas cost and deterministic $\Delta TCO$}
\label{tab:import_scenarios_tco_det}
\resizebox{\textwidth}{!}{
\begin{tabular}{llrr}
\toprule
\textbf{Scenario} 
& \textbf{Stratum} 
& $TCO_{\mathrm{BAU,gas}}$ (USD/user) 
& $\Delta TCO_{\mathrm{det}}$ (USD/user) \\
\midrule
50\% imported gas & E1 & 983.44  & -590.80 \\
50\% imported gas & E2 & 1027.38 & -627.82 \\
50\% imported gas & E3 & 1382.18 & -958.14 \\
50\% imported gas & E4 & 1702.25 & -1230.11 \\
50\% imported gas & E5 & 2569.79 & -2018.54 \\
50\% imported gas & E6 & 4645.17 & -3970.30 \\
\midrule
100\% imported gas & E1 & 1287.45 & -894.81 \\
100\% imported gas & E2 & 1344.97 & -945.41 \\
100\% imported gas & E3 & 1809.45 & -1385.41 \\
100\% imported gas & E4 & 2228.46 & -1756.33 \\
100\% imported gas & E5 & 3364.19 & -2812.94 \\
100\% imported gas & E6 & 6081.13 & -5406.25 \\
\bottomrule
\end{tabular}
}
\end{table}

The MAC results in Table~\ref{tab:import_scenarios_macc_ci} remain negative across all strata and under both imported-gas scenarios. Under the 50\% imported-gas scenario, deterministic MAC values range from USD $-192.60$/tCO$_2$e in Stratum~1 to USD $-776.15$/tCO$_2$e in Stratum~6. In the stochastic evaluation, mean MAC values range from USD $-176.90$/tCO$_2$e to USD $-771.74$/tCO$_2$e, and the 95\% intervals remain entirely below zero for all strata.

Under the 100\% imported-gas scenario, the MAC values become even more negative. Deterministic MAC values range from USD $-291.71$/tCO$_2$e in Stratum~1 to USD $-1056.86$/tCO$_2$e in Stratum~6. The stochastic mean ranges from USD $-272.15$/tCO$_2$e in Stratum~1 to USD $-1054.91$/tCO$_2$e in Stratum~6, with all 95\% intervals remaining below zero. Relative to the baseline case, both imported-gas scenarios shift the MAC distribution downward because the BAU gas-cooking alternative becomes more expensive. This effect is stronger under full import dependence, where the stochastic mean MAC reaches USD $-1054.91$/tCO$_2$e in Stratum~6. Therefore, greater gas-import dependence increases the monetary savings per ton of CO$_2$e avoided.

\begin{table}[htbp]
\centering
\caption{MAC under imported-gas scenarios in Bogotá: deterministic MAC and stochastic MAC distribution}
\label{tab:import_scenarios_macc_ci}
\resizebox{\textwidth}{!}{
\begin{tabular}{llccc}
\toprule
\textbf{Scenario} 
& \textbf{Stratum}  
& $MAC_{\mathrm{det}}$ (USD/tCO$_2$e) 
& $\overline{MAC}_{\mathrm{stoch}}$ (USD/tCO$_2$e)
& 95\% CI (USD/tCO$_2$e) \\
\midrule
50\% imported gas & E1 & -192.60 & -176.90 & [-217.78,\,-137.71] \\
50\% imported gas & E2 & -234.46 & -219.46 & [-277.89,\,-165.72] \\
50\% imported gas & E3 & -443.28 & -559.63 & [-683.87,\,-445.36] \\
50\% imported gas & E4 & -462.10 & -581.40 & [-678.36,\,-492.70] \\
50\% imported gas & E5 & -602.74 & -739.45 & [-886.48,\,-604.21] \\
50\% imported gas & E6 & -776.15 & -771.74 & [-918.77,\,-636.50] \\
\midrule
100\% imported gas & E1 & -291.71 & -272.15 & [-323.53,\,-224.51] \\
100\% imported gas & E2 & -353.06 & -331.99 & [-410.16,\,-260.23] \\
100\% imported gas & E3 & -640.95 & -795.61 & [-960.20,\,-646.09] \\
100\% imported gas & E4 & -659.77 & -819.52 & [-942.80,\,-707.51] \\
100\% imported gas & E5 & -839.95 & -1022.62 & [-1219.35,\,-842.05] \\
100\% imported gas & E6 & -1056.86 & -1054.91 & [-1251.65,\,-874.35] \\
\bottomrule
\end{tabular}
}
\end{table}

%%%%%%%%%%%%%%%%%%%%%%%%%%%%%%%%%%%%%%%%%%%%%%%%%%%%%%%%%%
%%%%%%%%%%%%%%%%%%%%%%%%%%%%%%%%%%%%%%%%%%%%%%%%%%%%%%%%%%

\bibliographystyle{elsarticle-harv}
\bibliography{ecooking_ref}

%%%%%%%%%%%%%%%%%%%%%%%%%%%%%%%%%%%%%%%%%%%%%%%%%%%%%%%%
%%%%%%%%%%%%%%%%%%%%%%%%%%%%%%%%%%%%%%%%%%%%%%%%%%%%%%%%

\end{document}